% !TeX root = 03_intro_ml_polytech.tex
% !TeX encoding = UTF-8
% !TeX spellcheck = en_GB
\begin{frame}{Supervised Learning}{(Re-)introducing labels\dots}
    Differently from unsupervised learning, \textbf{supervised learning} deals with \textbf{labeled} datasets:
    \begin{equation*}
        \mathcal{D} = \left\lbrace \left(\vec{x}_i, y_i\right) \right\rbrace_{i = 1}^N
    \end{equation*}

    \vfill \pause

    The key point is to find:
    \begin{equation*}
        f\colon \mathbb{R}^{d_1 \times d_2 \times \dots} \to \mathbb{R}^{k_1 \times k_2 \times \dots}
    \end{equation*}
    {\footnotesize
    For instance, we can have:
    \begin{itemize}
        \item \textbf{tabular data} $X \in \mathbb{R}^{N \times P}$:
              \begin{itemize}
                  \item \textbf{regression:} ($\mathbb{R}$) \emph{find the critical heat flux value, given conditions in a nuclear plant}
                  \item \textbf{classification:} ($\mathbb{R}^{k_1 \times k_2 \times \dots}$) \emph{assign a mineral class to a series of spectroscopic values}
                  \item \dots
              \end{itemize}
        \item \textbf{images} $X \in \mathbb{R}^{C \times H \times W}$:
              \begin{itemize}
                  \item \textbf{regression:} ($\mathbb{R}$) \emph{find the Reynolds number given a thermofluidic slice}
                  \item \textbf{segmentation:} ($\mathbb{R}^{H \times W}$) \emph{isolate pixels of different animals/minerals/objects}
                  \item \dots
              \end{itemize}
        \item \textbf{videos} $X \in \mathbb{R}^{T \times C \times H \times W}$:
              \begin{itemize}
                  \item \textbf{segmentation:} ($\mathbb{R}^{T \times H \times W}$) \emph{follow pixels of different animals/minerals/objects}
                  \item \textbf{detection:} ($\mathbb{R}$) \emph{detect a phase transition point in a finite-element simulation}
                  \item \dots
              \end{itemize}
    \end{itemize}
    }
\end{frame}

\begin{frame}{Supervised Methods}{Building some intuition\dots}
    Let us start by something which can be \textbf{linearly separated}:
    \vfill
    \begin{columns}[T]
        \begin{column}{0.32\textwidth}
            \begin{center}
                \begin{tikzpicture}[
        ax/.style={CEAblack, thick, ->},
        classif/.style={CEAgreen, ultra thick, dashed}
    ]
    \pgfmathsetseed{42} % fixed seed -> reproducible figure
    \newcommand{\boxmuller}{%
        \pgfmathsetmacro{\bmR}{sqrt(-2*ln(1-rnd))}%
        \pgfmathsetmacro{\bmT}{360*rnd}% degrees: PGF trig is in degrees
        \pgfmathsetmacro{\gx}{\bmR*cos(\bmT)}%
        \pgfmathsetmacro{\gy}{\bmR*sin(\bmT)}%
    }

    % Axes
    \draw[ax] (-0.25, 0.0) -- (3.5, 0.0) node[below] {$x_1$};
    \draw[ax] (0.0, -0.25) -- (0.0, 3.0) node[left] {$x_2$};

    % Points
    \foreach \i in {1,...,25} {
            \boxmuller
            \pgfmathsetmacro{\px}{2.5 + 0.3*\gx}
            \pgfmathsetmacro{\py}{0.5 + 0.3*\gy}
            \fill[CEAred!70] (\px,\py) circle (1.5pt);
        }
    \foreach \i in {1,...,25} {
            \boxmuller
            \pgfmathsetmacro{\px}{0.5 + 0.3*\gx}
            \pgfmathsetmacro{\py}{2.5 + 0.4*\gy}
            \fill[CEAblue!70] (\px,\py) circle (1.5pt);
        }

    % Classification
    \draw[classif] (0.95, -0.1) -- (1.99, 2.9);
\end{tikzpicture}


                \textbf{Logistic regression} provides a \emph{thresholding} framework to classify objects
            \end{center}
        \end{column}
        \hfil
        \begin{column}<2->{0.32\textwidth}
            \begin{center}
                \begin{tikzpicture}[
        ax/.style={CEAblack, thick, ->},
        classif/.style={CEAgreen, ultra thick, dashed},
        corridor/.style={CEAgreen, thick, dotted},
        dist/.style={CEAgreen, thick, <->}
    ]
    \pgfmathsetseed{42} % fixed seed -> reproducible figure
    \newcommand{\boxmuller}{%
        \pgfmathsetmacro{\bmR}{sqrt(-2*ln(1-rnd))}%
        \pgfmathsetmacro{\bmT}{360*rnd}% degrees: PGF trig is in degrees
        \pgfmathsetmacro{\gx}{\bmR*cos(\bmT)}%
        \pgfmathsetmacro{\gy}{\bmR*sin(\bmT)}%
    }

    % Axes
    \draw[ax] (-0.25, 0.0) -- (3.5, 0.0) node[below] {$x_1$};
    \draw[ax] (0.0, -0.25) -- (0.0, 3.0) node[left] {$x_2$};

    % Points
    \foreach \i in {1,...,25} {
            \boxmuller
            \pgfmathsetmacro{\px}{2.5 + 0.3*\gx}
            \pgfmathsetmacro{\py}{0.5 + 0.3*\gy}
            \fill[CEAred!70] (\px,\py) circle (1.5pt);
        }
    \foreach \i in {1,...,25} {
            \boxmuller
            \pgfmathsetmacro{\px}{0.5 + 0.3*\gx}
            \pgfmathsetmacro{\py}{2.5 + 0.4*\gy}
            \fill[CEAblue!70] (\px,\py) circle (1.5pt);
        }

    % Classification
    \draw[classif] (0.05, -0.1) -- (2.85, 2.9);
    \draw[corridor] (1.05, -0.25) -- (2.85, 1.7) node[right] {$\ell_+$};
    \draw[corridor] (0.05, 0.9) -- (1.95, 3.0) node[right] {$\ell_-$};
    \draw[dist] (1.2, 2.1) -- (2.25, 1.05);
\end{tikzpicture}


                \only<2-3>{%
                    What if we asked to create the \enquote{most separation} between the objects?
                }
            \end{center}
        \end{column}
        \hfil
        \begin{column}<3->{0.32\textwidth}
            \begin{center}
                \begin{tikzpicture}[
        ax/.style={CEAblack, thick, ->},
        classif/.style={CEAgreen, ultra thick, dashed},
    ]
    \pgfmathsetseed{42} % fixed seed -> reproducible figure
    \newcommand{\boxmuller}{%
        \pgfmathsetmacro{\bmR}{sqrt(-2*ln(1-rnd))}%
        \pgfmathsetmacro{\bmT}{360*rnd}% degrees: PGF trig is in degrees
        \pgfmathsetmacro{\gx}{\bmR*cos(\bmT)}%
        \pgfmathsetmacro{\gy}{\bmR*sin(\bmT)}%
    }

    % Axes
    \draw[ax] (-0.25, 0.0) -- (3.5, 0.0) node[below] {$x_1$};
    \draw[ax] (0.0, -0.25) -- (0.0, 3.0) node[left] {$x_2$};

    % Points
    \foreach \i in {1,...,25} {
            \boxmuller
            \pgfmathsetmacro{\px}{2.5 + 0.3*\gx}
            \pgfmathsetmacro{\py}{0.5 + 0.3*\gy}
            \fill[CEAred!70] (\px,\py) circle (1.5pt);
        }
    \foreach \i in {1,...,25} {
            \boxmuller
            \pgfmathsetmacro{\px}{0.5 + 0.3*\gx}
            \pgfmathsetmacro{\py}{2.5 + 0.4*\gy}
            \fill[CEAblue!70] (\px,\py) circle (1.5pt);
        }

    % Classification
    \draw[classif]
    (0.05, -0.1) -- (0.05, 0.8) --
    (1.1, 0.8) -- (1.1, 1.2) --
    (1.9, 1.2) -- (1.9, 1.7) --
    (2.5, 1.7) -- (2.5, 2.4) --
    (2.9, 2.4) -- (2.9, 2.9);
\end{tikzpicture}


                \only<3>{%
                    What if we asked to \enquote{divide the space} in patches and assign the objects?
                }
            \end{center}
        \end{column}
    \end{columns}
\end{frame}

\begin{frame}{Support Vector Machines}{Optimising the corridor\dots}
    Let us define the two \textbf{corridor} lines:
    \begin{equation*}
        \ell_{\pm}\colon \vec{a} \cdot \vec{x} + b =
        \begin{cases}
            +1 & \qif y = 1~\text{(blue)} \\
            -1 & \qif y = -1~\text{(red)}
        \end{cases}
    \end{equation*}
    \vfill
    \begin{columns}[c]
        \begin{column}{0.4\textwidth}
            \centering
            \begin{tikzpicture}[
        ax/.style={CEAblack, thick, ->},
        classif/.style={CEAgreen, ultra thick, dashed},
        corridor/.style={CEAgreen, thick, dotted},
        dist/.style={CEAgreen, thick, <->}
    ]
    \pgfmathsetseed{42} % fixed seed -> reproducible figure
    \newcommand{\boxmuller}{%
        \pgfmathsetmacro{\bmR}{sqrt(-2*ln(1-rnd))}%
        \pgfmathsetmacro{\bmT}{360*rnd}% degrees: PGF trig is in degrees
        \pgfmathsetmacro{\gx}{\bmR*cos(\bmT)}%
        \pgfmathsetmacro{\gy}{\bmR*sin(\bmT)}%
    }

    % Axes
    \draw[ax] (-0.25, 0.0) -- (3.5, 0.0) node[below] {$x_1$};
    \draw[ax] (0.0, -0.25) -- (0.0, 3.0) node[left] {$x_2$};

    % Points
    \foreach \i in {1,...,25} {
            \boxmuller
            \pgfmathsetmacro{\px}{2.5 + 0.3*\gx}
            \pgfmathsetmacro{\py}{0.5 + 0.3*\gy}
            \fill[CEAred!70] (\px,\py) circle (1.5pt);
        }
    \foreach \i in {1,...,25} {
            \boxmuller
            \pgfmathsetmacro{\px}{0.5 + 0.3*\gx}
            \pgfmathsetmacro{\py}{2.5 + 0.4*\gy}
            \fill[CEAblue!70] (\px,\py) circle (1.5pt);
        }

    % Classification
    \draw[classif] (0.05, -0.1) -- (2.85, 2.9);
    \draw[corridor] (1.05, -0.25) -- (2.85, 1.7) node[right] {$\ell_+$};
    \draw[corridor] (0.05, 0.9) -- (1.95, 3.0) node[right] {$\ell_-$};
    \draw[dist] (1.2, 2.1) -- (2.25, 1.05);
\end{tikzpicture}

        \end{column}
        \hfil
        \begin{column}{0.6\textwidth}
            We can now \textbf{maximise} the distance between the lines $\ell_+$ and $\ell_-$:
            \begin{equation*}
                \operatorname{dist}(\ell_+, \ell_-)
                =
                \frac{2}{\norm{\vec{a}}}
            \end{equation*}
            constrained to:
            \begin{equation*}
                y_i \left( \vec{a} \cdot \vec{x}_i + b \right) \ge 1
            \end{equation*}
        \end{column}
    \end{columns}
\end{frame}

\begin{frame}{Support Vector Machines}{Optimising the corridor\dots}
    The condition becomes:
    \begin{equation*}
        \left(\vec{a}^*, b^*\right)
        =
        \arg\min\limits_{\vec{a}, b}
        \frac{\norm{\vec{a}}}{2}
        +
        \only<1>{\sum\limits_{i = 1}^N}
        \only<2>{\alert{\sum\limits_{i = 1}^S}}
        \mu_i y_i \left( \vec{a} \cdot \vec{x}_i + b\right)
    \end{equation*}
    \vfill
    \begin{columns}[c]
        \begin{column}{0.4\textwidth}
            \centering
            \begin{tikzpicture}[
        ax/.style={CEAblack, thick, ->},
        classif/.style={CEAgreen, ultra thick, dashed},
        corridor/.style={CEAgreen, thick, dotted},
        dist/.style={CEAgreen, thick, <->}
    ]
    \pgfmathsetseed{42} % fixed seed -> reproducible figure
    \newcommand{\boxmuller}{%
        \pgfmathsetmacro{\bmR}{sqrt(-2*ln(1-rnd))}%
        \pgfmathsetmacro{\bmT}{360*rnd}% degrees: PGF trig is in degrees
        \pgfmathsetmacro{\gx}{\bmR*cos(\bmT)}%
        \pgfmathsetmacro{\gy}{\bmR*sin(\bmT)}%
    }

    % Axes
    \draw[ax] (-0.25, 0.0) -- (3.5, 0.0) node[below] {$x_1$};
    \draw[ax] (0.0, -0.25) -- (0.0, 3.0) node[left] {$x_2$};

    % Points
    \foreach \i in {1,...,25} {
            \boxmuller
            \pgfmathsetmacro{\px}{2.5 + 0.3*\gx}
            \pgfmathsetmacro{\py}{0.5 + 0.3*\gy}
            \fill[CEAred!70] (\px,\py) circle (1.5pt);
        }
    \foreach \i in {1,...,25} {
            \boxmuller
            \pgfmathsetmacro{\px}{0.5 + 0.3*\gx}
            \pgfmathsetmacro{\py}{2.5 + 0.4*\gy}
            \fill[CEAblue!70] (\px,\py) circle (1.5pt);
        }

    % Classification
    \draw[classif] (0.05, -0.1) -- (2.85, 2.9);
    \draw[corridor] (1.05, -0.25) -- (2.85, 1.7) node[right] {$\ell_+$};
    \draw[corridor] (0.05, 0.9) -- (1.95, 3.0) node[right] {$\ell_-$};
    \draw[dist] (1.2, 2.1) -- (2.25, 1.05);
\end{tikzpicture}

        \end{column}
        \hfil
        \begin{column}<2->{0.6\textwidth}
            But we know that the \textbf{constraint} is active \textbf{only} for the points on the \textbf{border}:

            \begin{center}
                \large
                The method only depends on $S$ \textbf{support vectors} and is hence called:

                \vspace{0.5em}

                {\huge
                    Support Vector Machine
                }
            \end{center}

            \begin{equation*}
                \vec{a}^{*}
                =
                \sum\limits_{i = 1}^S \mu_i y_i \vec{x}_i,
                \quad
                b^*
                =
                y_i - \vec{a}^* \cdot \vec{x}_{\hat\i}
            \end{equation*}
        \end{column}
    \end{columns}
\end{frame}

\begin{frame}{Support Vector Machines}{Creating non-linear separation\dots}
    The only issue so far is that we worked with linearly separarable samples.
    But what if?
    \vfill
    \begin{center}
        \begin{tikzpicture}
    % Axes
    \draw[thick, CEAblack, ->] (-2.0, 0.0) -- (2.0, 0.0) node[below] {$x_1$};
    \draw[thick, CEAblack, ->] (0.0, -2.0) -- (0.0, 2.0) node[left] {$x_2$};

    \pgfmathsetseed{42} % fixed seed -> reproducible figure

    % Box–Muller: after \boxmuller, \gx and \gy are two independent N(0,1) draws
    \newcommand{\boxmuller}{%
        \pgfmathsetmacro{\bmR}{sqrt(-2*ln(1-rnd))}%
        \pgfmathsetmacro{\bmT}{360*rnd}% degrees: PGF trig is in degrees
        \pgfmathsetmacro{\gx}{\bmR*cos(\bmT)}%
        \pgfmathsetmacro{\gy}{\bmR*sin(\bmT)}%
    }

    \foreach \i in {1,...,100} {
            \boxmuller
            \pgfmathsetmacro{\th}{360*rnd}
            \pgfmathsetmacro{\rr}{1.5 + 0.25*\gx}
            \pgfmathsetmacro{\px}{0.0 + \rr*cos(\th)}
            \pgfmathsetmacro{\py}{0.0 + \rr*sin(\th)}
            \fill[CEAred!70] (\px,\py) circle (1.5pt);
        }

    \foreach \i in {1,...,100} {
            \boxmuller
            \pgfmathsetmacro{\th}{360*rnd}
            \pgfmathsetmacro{\rr}{0.25 + 0.06*\gx}
            \pgfmathsetmacro{\px}{0.0 + \rr*cos(\th)}
            \pgfmathsetmacro{\py}{0.0 + \rr*sin(\th)}
            \fill[CEAblue!70] (\px,\py) circle (1.5pt);
        }
\end{tikzpicture}

    \end{center}
\end{frame}

\begin{frame}{Support Vector Machines}{Creating non-linear separation\dots}
    The idea is \enquote{artificially} create the separation by \enquote{adding dimensions}:
    \begin{equation*}
        \vec{x} = (x_1, x_2) \in \mathbb{R}^2
        \quad \longmapsto \quad
        \vec{x}^{\prime} = (x_1, x_2, e^{-(x_1^2 + x_2^2)}) \in \mathbb{R}^3
    \end{equation*}

    \vfill

    \begin{center}
        \begin{tikzpicture}
    % Axes
    \draw[thick, CEAblack, ->] (-2.0, 0.0) -- (2.0, 0.0) node[below] {$x_1$};
    \draw[thick, CEAblack, ->] (0.0, -2.0) -- (0.0, 2.0) node[left] {$x_2$};

    \pgfmathsetseed{42} % fixed seed -> reproducible figure

    % Box–Muller: after \boxmuller, \gx and \gy are two independent N(0,1) draws
    \newcommand{\boxmuller}{%
        \pgfmathsetmacro{\bmR}{sqrt(-2*ln(1-rnd))}%
        \pgfmathsetmacro{\bmT}{360*rnd}% degrees: PGF trig is in degrees
        \pgfmathsetmacro{\gx}{\bmR*cos(\bmT)}%
        \pgfmathsetmacro{\gy}{\bmR*sin(\bmT)}%
    }

    \foreach \i in {1,...,100} {
            \boxmuller
            \pgfmathsetmacro{\th}{360*rnd}
            \pgfmathsetmacro{\rr}{1.5 + 0.25*\gx}
            \pgfmathsetmacro{\px}{0.0 + \rr*cos(\th)}
            \pgfmathsetmacro{\py}{0.0 + \rr*sin(\th)}
            \fill[CEAred!70] (\px,\py) circle (1.5pt);
        }

    \foreach \i in {1,...,100} {
            \boxmuller
            \pgfmathsetmacro{\th}{360*rnd}
            \pgfmathsetmacro{\rr}{0.25 + 0.06*\gx}
            \pgfmathsetmacro{\px}{0.0 + \rr*cos(\th)}
            \pgfmathsetmacro{\py}{0.0 + \rr*sin(\th)}
            \fill[CEAblue!70] (\px,\py) circle (1.5pt);
        }

    % Arrow
    \begin{scope}[xshift=3cm]
        \node at (0, 0) {\Huge $\Rightarrow$};
    \end{scope}

    % 3D plot
    \begin{scope}[xshift=5cm]
        \draw[thick, CEAblack, ->] (0.0, 0.0) -- (2.0, 0.0) node[below] {$x_1$};
        \draw[thick, CEAblack, ->] (0.0, 0.0) -- (-1.0, -1.0) node[left] {$x_2$};
        \draw[thick, CEAblack, ->] (0.0, 0.0) -- (0.0, 2.0) node[left] {$x_3$};

        \foreach \n in {0,...,50} {
                \pgfmathsetmacro{\px}{2.0 * rnd - 1.0};
                \pgfmathsetmacro{\py}{0.5 * (rnd - 0.65)};
                \fill[CEAred!70] (\px,\py) circle (1.5pt);
                \pgfmathsetmacro{\qx}{2.0 * rnd - 1.0};
                \pgfmathsetmacro{\qy}{0.5 * (rnd - 0.25)};
                \fill[CEAred!70, opacity=0.5] (\qx,\qy) circle (1.5pt);
                \pgfmathsetmacro{\ux}{0.5 * (rnd - 0.5)};
                \pgfmathsetmacro{\uy}{0.5 * (rnd + 2.5)};
                \fill[CEAblue!70] (\ux,\uy) circle (1.5pt);
                \pgfmathsetmacro{\vx}{0.5 * (rnd - 0.5)};
                \pgfmathsetmacro{\vy}{0.5 * (rnd + 2.5)};
                \fill[CEAblue!70, opacity=0.5] (\vx,\vy) circle (1.5pt);
            }

        \draw[CEAblack!70, fill=CEAblack!15, fill opacity=0.25] (-1.5, 0.7) -- (1.5, 0.7) -- (1.9, 1.1) -- (-1.1, 1.1) -- cycle;
    \end{scope}
\end{tikzpicture}

    \end{center}
\end{frame}

\begin{frame}{Support Vector Machines}{Creating non-linear separation\dots}
    This is called \textbf{kernel trick}:
    \begin{equation*}
        \vec{x} \mapsto \vec{\phi}(\vec{x})
        \quad \Rightarrow \quad
        \vec{a}^*
        =
        \sum\limits_{i = 1}^S \mu_i y_i \vec{\phi}(\vec{x}_i)
    \end{equation*}
    which means that the classification function becomes:
    \begin{equation*}
        f(\vec{x}_j)
        =
        \vec{a}^* \cdot \vec{x}_j + b
        \quad \longrightarrow \quad
        f(\vec{x}_j)
        =
        \sum\limits_{i = 1}^S \mu_i y_i \underbrace{\vec{\phi}(\vec{x}_i) \cdot \vec{\phi}(\vec{x}_j)}_{= K(\vec{x}_i, \vec{x}_j)} + b
    \end{equation*}

    \vfill \pause

    We only need to know the combination:
    \begin{equation*}
        K(\vec{x}_i, \vec{x}_j)
        =
        \begin{cases}
            \vec{x}_i \cdot \vec{x}_j                                & \text{linear kernel}                      \\
            (1 + \vec{x}_i \cdot \vec{x}_j)^p                        & \text{polynomial kernel}                  \\
            \exp(-\frac{\norm{\vec{x}_i - \vec{x}_j}^2}{2 \sigma})^p & \text{RBF (\emph{radial basis function})} \\
            \dots
        \end{cases}
    \end{equation*}
\end{frame}

\begin{frame}{Decision Trees}{Separating into buckets\dots}
    Another useful \textbf{supervised} algorithm uses \textbf{partitions} of the space to \enquote{assign samples to buckets}:

    \vfill

    \begin{columns}[c]
        \begin{column}{0.4\textwidth}
            \centering
            \begin{tikzpicture}[
        ax/.style={CEAblack, thick, ->},
        classif/.style={CEAgreen, ultra thick, dashed},
    ]
    \pgfmathsetseed{42} % fixed seed -> reproducible figure
    \newcommand{\boxmuller}{%
        \pgfmathsetmacro{\bmR}{sqrt(-2*ln(1-rnd))}%
        \pgfmathsetmacro{\bmT}{360*rnd}% degrees: PGF trig is in degrees
        \pgfmathsetmacro{\gx}{\bmR*cos(\bmT)}%
        \pgfmathsetmacro{\gy}{\bmR*sin(\bmT)}%
    }

    % Axes
    \draw[ax] (-0.25, 0.0) -- (3.5, 0.0) node[below] {$x_1$};
    \draw[ax] (0.0, -0.25) -- (0.0, 3.0) node[left] {$x_2$};

    % Points
    \foreach \i in {1,...,25} {
            \boxmuller
            \pgfmathsetmacro{\px}{2.5 + 0.3*\gx}
            \pgfmathsetmacro{\py}{0.5 + 0.3*\gy}
            \fill[CEAred!70] (\px,\py) circle (1.5pt);
        }
    \foreach \i in {1,...,25} {
            \boxmuller
            \pgfmathsetmacro{\px}{0.5 + 0.3*\gx}
            \pgfmathsetmacro{\py}{2.5 + 0.4*\gy}
            \fill[CEAblue!70] (\px,\py) circle (1.5pt);
        }

    % Classification
    \draw[classif]
    (0.05, -0.1) -- (0.05, 0.8) --
    (1.1, 0.8) -- (1.1, 1.2) --
    (1.9, 1.2) -- (1.9, 1.7) --
    (2.5, 1.7) -- (2.5, 2.4) --
    (2.9, 2.4) -- (2.9, 2.9);
\end{tikzpicture}

        \end{column}
        \hfil
        \begin{column}{0.6\textwidth}
            \centering
            \begin{tikzpicture}[
        scale=0.65,
        every node/.style={scale=0.65},
        leaf/.style={
                circle,
                CEAblack,
                fill=CEAblack!15,
                align=center,
                drop shadow = {shadow xshift = 2pt, shadow yshift = -2pt,
                        opacity = 0.35, fill = black},
            }
    ]
    \node<1>[set] (n0) at (0, 0) {question};
    \node<1>[set, yshift=-1em, below left=0.35em of n0] (n00) {question};
    \node<1>[set, yshift=-1em, below right=0.35em of n0] (n01) {question};
    \node<1>[set, yshift=-1em, below left=0.35em of n00] (n000) {question};
    \node<1>[leaf, yshift=-1em, below right=0.15em of n00] (n001) {decision};
    \node<1>[set, yshift=-1em, below left=0.15em of n000] (n0000) {question};
    \node<1>[set, yshift=-1em, below right=0.35em of n000] (n0001) {question};
    \node<1>[leaf, yshift=-1em, below left=0.35em of n0000] (n00000) {decision};
    \node<1>[leaf, yshift=-1em, below right=0.35em of n0000] (n00001) {decision};

    \node<2->[set] (n0) at (0, 0) {$x_1 < 0.3$};
    \node<2->[set, yshift=-1em, below left=0.35em of n0] (n00) {$x_2 > 0.1$};
    \node<2->[set, yshift=-1em, below right=0.35em of n0] (n01) {$x_2 < 0.7$};
    \node<2->[set, yshift=-1em, below left=0.35em of n00] (n000) {$x_1 > 1.5$};
    \node<2->[leaf, yshift=-1em, below right=0.15em of n00] (n001) {class = 1};
    \node<2->[set, yshift=-1em, below left=0.15em of n000] (n0000) {$x_2 < 0.5$};
    \node<2->[set, yshift=-1em, below right=0.35em of n000] (n0001) {$x_1 < -0.1$};
    \node<2->[leaf, yshift=-1em, below left=0.35em of n0000] (n00000) {class = 0};
    \node<2->[leaf, yshift=-1em, below right=0.35em of n0000] (n00001) {class = 2};

    \draw[->, CEAblack] (n0.south) -- (n00.north east);
    \draw[->, CEAblack] (n0.south) -- (n01.north west);
    \draw[->, CEAblack] (n00.south) -- (n000.north east);
    \draw[->, CEAblack] (n00.south) -- (n001.north west);
    \draw[->, CEAblack] (n000.south) -- (n0000.north east);
    \draw[->, CEAblack] (n000.south) -- (n0001.north west);
    \draw[->, CEAblack] (n0000.south) -- (n00000.north east);
    \draw[->, CEAblack] (n0000.south) -- (n00001.north west);
\end{tikzpicture}

        \end{column}
    \end{columns}
\end{frame}

\begin{frame}{Decision Trees}{Separating into buckets\dots}
    Like SVMs, \textbf{decision trees} are trained via \textbf{gradient descent} against particular losses:
    \begin{itemize}
        \item \textbf{SVM for classification:} standard cross-entropy
        \item \textbf{SVM for regression:} standard MSE
        \item \textbf{Decision trees for classification:} \emph{Gini impurity} or cross-entropy
        \item \textbf{Decision trees for regression:} standard MSE
    \end{itemize}

    \vfill \pause

    The main issue with trees is that
    \begin{center}
        \emph{the natural outcome \enquote{without stopping} is \enquote{one point per bucket/leaf}}
    \end{center}

    \vfill \pause

    \includegraphics[height=12pt]{01_slides/img/brain.pdf}~Are they a high-variance or a high-bias algorithm?

    \vfill \pause

    \begin{center}
        \textbf{Decision trees} are very high \textbf{variance} algorithms (they tend to overfit)
    \end{center}
\end{frame}

\begin{frame}{Decision Trees}{Separating into buckets\dots}
    To mitigate:

    \vfill

    \begin{itemize}
        \item \textbf{Pruning:} create some rules to stop from getting to the end
              \begin{itemize}
                  \item Limit the number of \textbf{leaf nodes} (terminal nodes)
                  \item Randomise the \textbf{input features}
                  \item Limit the \textbf{depth level}
                  \item \dots
              \end{itemize}
        \item \textbf{Ensembling:} train multiple pruned decision trees and take the average value $\Rightarrow$ these are called \textbf{Random Forests}
        \item \textbf{Boosting:} iteratively train a tree and re-weight the badly classified/regresses points $\Rightarrow$ the procedure is called \textbf{Gradient Boosting}
    \end{itemize}

    \vfill
\end{frame}

\begin{frame}[fragile]{Supervised Learning}{\includegraphics[height=12pt]{01_slides/img/tool.pdf}~Working on spectroscopy data}
    \begin{itemize}
        \item \alert<2>{\includegraphics[height=12pt]{01_slides/img/pen-tool.pdf}~Open the \texttt{machine\_learning.ipynb} notebook and run all cells}
        \item \alert<3>{\includegraphics[height=12pt]{01_slides/img/pen-tool.pdf}~Import the \textbf{Support Vector Classifier} class (\texttt{SVC}) from Scikit-learn}
        \item \alert<4>{\includegraphics[height=12pt]{01_slides/img/pen-tool.pdf}~Train the \texttt{SVC} on the PCA-transformed data, then predict the labels of training and validation data}
        \item \alert<5-6>{\includegraphics[height=12pt]{01_slides/img/pen-tool.pdf}~Compute the confusion matrix, the accuracy, the precision, the recall, and the F1 score}
        \item \alert<7>{\includegraphics[height=12pt]{01_slides/img/pen-tool.pdf}~Import the \texttt{XGBClassifier} from \texttt{xgboost}}
        \item \alert<8>{\includegraphics[height=12pt]{01_slides/img/pen-tool.pdf}~Train the gradient boosting algorithm on the PCA-transformed data, then predict the labels of training and validation data}
        \item \alert<9-10>{\includegraphics[height=12pt]{01_slides/img/pen-tool.pdf}~Compute the confusion matrix, the accuracy, the precision, the recall, and the F1 score}
        \item \alert<11>{\includegraphics[height=12pt]{01_slides/img/brain.pdf}~What can we say about the two classification algorithms?}
    \end{itemize}

    \begin{onlyenv}<3>
        \begin{tikzpicture}[overlay]
            \node[set, text width=30em, xshift=-2em, yshift=-6em] at (current page.center) {
                \begin{notebookin}
from sklearn.svm import SVC
                \end{notebookin}
            };
        \end{tikzpicture}
    \end{onlyenv}

    \begin{onlyenv}<4>
        \begin{tikzpicture}[overlay]
            \node[set, text width=38em, xshift=-2em, yshift=-6em] at (current page.center) {
                \begin{notebookin}
svc = SVC(kernel="rbf")
svc.fit(X_train_pca, y_train_sel-1)  # use -1 to start from zero
                \end{notebookin}
                \begin{notebookin}
y_train_pred = svc.predict(X_train_pca)
y_val_pred = svc.predict(X_val_pca)
                \end{notebookin}
            };
        \end{tikzpicture}
    \end{onlyenv}

    \begin{onlyenv}<5>
        \begin{tikzpicture}[overlay]
            \node[set, text width=41em, xshift=-3em, yshift=-6em] at (current page.center) {
                \begin{notebookin}
from sklearn.metrics import confusion_matrix, accuracy_score, precision_score, recall_score, f1_score

for name, true, pred in [("train", y_train_sel, y_train_pred),
                         ("val", y_val_sel, y_val_pred)]:
    acc = accuracy_score(true-1, pred); print(f"Accuracy ({name}) = {acc:.1%}")
    for score in [precision_score, recall_score, f1_score]:
        value = score(true-1, pred, average="macro")
        print(f"{score.__name__.capitalize().removesuffix('_score')} ({name}) = {value:.1%}")
                \end{notebookin}
                \begin{notebookout}
Accuracy (train) = 73.5%
Precision (train) = 56.8%
Recall (train) = 45.8%
F1 (train) = 44.5%
Accuracy (val) = 72.2%
Precision (val) = 52.4%
Recall (val) = 46.1%
F1 (val) = 44.7%
                \end{notebookout}
            };
        \end{tikzpicture}
    \end{onlyenv}

    \begin{onlyenv}<6>
        \begin{tikzpicture}[overlay]
            \node[set, text width=38em, xshift=-3em, yshift=-6em] at (current page.center) {
                \centering
                \includegraphics[width=0.45\textwidth]{03_slides/img/svc_conf_train.png}
                \includegraphics[width=0.45\textwidth]{03_slides/img/svc_conf_val.png}
            };
        \end{tikzpicture}
    \end{onlyenv}

    \begin{onlyenv}<7>
        \begin{tikzpicture}[overlay]
            \node[set, text width=30em, xshift=-2em, yshift=-6em] at (current page.center) {
                \begin{notebookin}
from xgboost import XGBClassifier
                \end{notebookin}
            };
        \end{tikzpicture}
    \end{onlyenv}

    \begin{onlyenv}<8>
        \begin{tikzpicture}[overlay]
            \node[set, text width=37em, xshift=-2em, yshift=-6em] at (current page.center) {
                \begin{notebookin}
boost = XGBClassifier(n_estimators=250)
boost.fit(X_train_pca, y_train_sel-1)  # use -1 to start from zero
                \end{notebookin}
                \begin{notebookin}
y_train_pred = boost.predict(X_train_pca)
y_val_pred = boost.predict(X_val_pca)
                \end{notebookin}
            };
        \end{tikzpicture}
    \end{onlyenv}

    \begin{onlyenv}<9>
        \begin{tikzpicture}[overlay]
            \node[set, text width=41em, xshift=-3em, yshift=-6em] at (current page.center) {
                \begin{notebookin}
# same as before
                \end{notebookin}
                \begin{notebookout}
Accuracy (train) = 83.5%
Precision (train) = 82.9%
Recall (train) = 68.0%
F1 (train) = 72.5%
Accuracy (val) = 71.4%
Precision (val) = 51.8%
Recall (val) = 47.5%
F1 (val) = 47.9%
                \end{notebookout}
            };
        \end{tikzpicture}
    \end{onlyenv}

    \begin{onlyenv}<10>
        \begin{tikzpicture}[overlay]
            \node[set, text width=38em, xshift=-3em, yshift=-6em] at (current page.center) {
                \centering
                \includegraphics[width=0.45\textwidth]{03_slides/img/xgb_conf_train.png}
                \includegraphics[width=0.45\textwidth]{03_slides/img/xgb_conf_val.png}
            };
        \end{tikzpicture}
    \end{onlyenv}
\end{frame}

\begin{frame}{Supervised and Unsupervised ML}{Preparing for next time\dots}
    You should now be proficient in:
    \begin{itemize}
        \item Unsupervised ML methods for data exploration and visualisation
        \item Non linear supervised ML methods for classification and regression
    \end{itemize}
    \vfill \pause
    \begin{center}
        \huge
        Remember where you saved the data! Next time we shall start from there!
        \\[1em]
        Are there questions?
    \end{center}
    \vfill
\end{frame}
